\begin{ouverture}[college1]{4}{Solutions, solubilité et pH}{Partie 1 · Organisation et transformations de la matière}{Une eau limpide est-elle forcément pure ?}
\objectif{Distinguer dissolution, miscibilité et séparation.}
\objectif{Estimer une solubilité à partir de données.}
\objectif{Mesurer un pH et reconnaître un caractère acide ou basique.}
\prerequis{Mélange homogène et unités de masse et de volume.}
\end{ouverture}
\pcPage{Le cours}{Solutions, solubilité et pH}
\phantomsection\label{ess-4}
\essSection{1}{Une solution est un mélange}
Le solvant dissout une ou plusieurs espèces appelées solutés. Dans une solution aqueuse, le solvant est l’eau. Dissoudre du sel ne le détruit pas : sa présence peut être mise en évidence autrement. Deux liquides sont miscibles s’ils forment un mélange homogène. Filtrer retient des particules solides assez grosses, mais ne retire pas les espèces dissoutes.\par
\pcFigure{figures/c04-dissolution.pdf}{Une dissolution peut donner un mélange transparent contenant plusieurs espèces.}
\essSection{2}{Il existe une limite}
À température donnée, la solubilité indique la quantité maximale de soluté que l’on peut dissoudre dans une quantité de solvant définie. On précise ici les masses solubles pour $100\unite{mL}$ d’eau. Au-delà, une partie du solide reste présente : la solution est saturée. Agiter accélère souvent la dissolution, mais ne permet pas de dépasser indéfiniment cette limite. Les gaz peuvent aussi se dissoudre dans l’eau.\par
\clearpage
\essSection{3}{Lire une échelle de pH}
Pour les solutions aqueuses usuelles étudiées vers $25\,^{\circ}\mathrm C$, un pH inférieur à 7 indique un caractère acide ; un pH égal à 7, la neutralité ; un pH supérieur à 7, un caractère basique. Le papier pH donne une estimation ; le pH-mètre fournit une mesure à contrôler et étalonner. Le pH est sans unité. La couleur d’une solution n’indique pas directement son pH.\par
\pcFigure{figures/c04-ph.pdf}{L’échelle usuelle de 0 à 14 est un repère scolaire pour des solutions aqueuses courantes ; elle ne couvre pas toutes les situations possibles.}
\essSection{4}{Préserver l’eau et préparer la suite}
Une filtration ne rend pas à elle seule une eau potable. La distillation sépare notamment l’eau de certains solutés non volatils, sans constituer une garantie universelle contre tout contaminant. Le dioxyde de carbone dissous contribue à l’acidification de l’eau ; une baisse du pH n’implique pas nécessairement de passer sous 7. \textbf{Repère pour la 3e :} les ions $\ce{H+}$ et $\ce{OH-}$ participent au caractère acide ou basique ; leur modèle sera approfondi, sans calcul ionique ici.\par
\begin{pcEncadre}{method}{Un exemple expliqué}À une température donnée, on fournit une limite de $36\unite g$ de sel pour $100\unite{mL}$ d’eau. Dans $50\unite{mL}$, la masse maximale vaut $36\times50/100=18\unite g$. Si l’on ajoute $22\unite g$, il reste $22-18=4\unite g$ non dissous, après agitation suffisante. Ce calcul utilise la quantité initiale d’eau, pas le volume final de solution.\end{pcEncadre}
\pcPage{Expérimenter et retenir}{Solutions, solubilité et pH}
\phantomsection\label{activite-4}
\begin{pcEncadre}{method}{Investigation en classe}\textbf{Question.} Comment estimer une limite de dissolution ? Avec les quantités et produits choisis par le professeur, ajouter de petites masses connues de sel dans un volume fixé d’eau à température stable. Agiter assez longtemps entre deux ajouts. Noter la dernière masse entièrement dissoute et la première laissant un dépôt persistant. En déduire un encadrement plutôt qu’une valeur artificiellement exacte. Sur d’autres échantillons autorisés, estimer le pH avec des bandelettes sans contaminer les flacons. Aucun mélange de produits ménagers.\end{pcEncadre}
\pcSousTitre{Ce que je dois savoir expliquer}\begin{questions}
\item Une solution contient un solvant et un ou plusieurs solutés.
\item La solubilité dépend du soluté, du solvant et de la température.
\item Une filtration ordinaire ne retire pas le sel dissous.
\item Acide, neutre et basique se distinguent par une mesure de pH.
\item Une eau transparente n’est pas nécessairement pure ni potable.
\end{questions}
\begin{pcEncadre}{attention}{Une erreur à éviter}« pH 6 est deux fois moins acide que pH 3. » Le pH n’est pas une échelle de proportionnalité ; on ne compare pas l’acidité en divisant les nombres de pH.\end{pcEncadre}
\pcRepere{cor-activite-4}{Pistes pour l’investigation}\pcGrille
\clearpage\markright{Je démarre}
\rubrique{Je démarre}
\par\noindent\begin{minipage}{\linewidth}
\begin{exercice}[1]{Le vocabulaire}
\begin{questions}
\item Dans une eau sucrée, nommer solvant et soluté.
\item Le mélange est-il pur ?
\end{questions}
\end{exercice}
\end{minipage}
\par\noindent\begin{minipage}{\linewidth}
\begin{exercice}[1]{Deux liquides}
\pcFigureEx{figures/ex-eau-huile.pdf}{}
\begin{questions}
\item Eau et huile forment deux couches après repos. Sont-elles miscibles ?
\item Comment nommer l’aspect du mélange ?
\end{questions}
\end{exercice}
\end{minipage}
\par\noindent\begin{minipage}{\linewidth}
\begin{exercice}[1]{Trois pH}
\begin{questions}
\item Classer les solutions de pH 3, 7 et 10 vers 25 °C.
\item Le pH porte-t-il une unité ?
\end{questions}
\end{exercice}
\end{minipage}
\par\noindent\begin{minipage}{\linewidth}
\begin{exercice}[1]{Un filtre en papier}
\begin{questions}
\item Peut-on retirer le sable d’une eau sableuse par filtration ?
\item Retirer le sel d’une eau salée par le même filtre ?
\end{questions}
\end{exercice}
\end{minipage}
\par\noindent\begin{minipage}{\linewidth}
\begin{exercice}[1]{Une solution saturée}
\begin{questions}
\item Après agitation prolongée, un dépôt de sel persiste. Quelle limite a probablement été atteinte ?
\item Pourquoi préciser la température ?
\end{questions}
\end{exercice}
\end{minipage}
\clearpage\markright{Je m’entraîne}
\rubrique{Je m’entraîne}
\par\noindent\begin{minipage}{\linewidth}
\begin{exercice}[2]{Adapter une quantité}
\begin{questions}
\item La limite fournie est $20\unite g$ pour $100\unite{mL}$ d’eau. Quelle masse maximale pour $250\unite{mL}$ ?
\item Tout le solide se dissout-il si l’on ajoute $60\unite g$ ?
\end{questions}
\end{exercice}
\end{minipage}
\par\noindent\begin{minipage}{\linewidth}
\begin{exercice}[2]{Conserver la masse}
\begin{questions}
\item On mélange $100\unite g$ d’eau et $8\unite g$ de sel, sans perte. Quelle masse totale obtient-on ?
\item Le volume total vaut-il forcément la somme des volumes initiaux ?
\end{questions}
\end{exercice}
\end{minipage}
\par\noindent\begin{minipage}{\linewidth}
\begin{exercice}[2]{Lire un encadrement}
\begin{questions}
\item $14\unite g$ se dissolvent complètement ; avec $16\unite g$, un dépôt persiste dans $50\unite{mL}$ d’eau. Encadrer la masse maximale.
\item Traduire pour $100\unite{mL}$ d’eau.
\end{questions}
\end{exercice}
\end{minipage}
\par\noindent\begin{minipage}{\linewidth}
\begin{exercice}[2]{Le gaz d’une boisson}
\begin{questions}
\item Pourquoi une boisson peut-elle former des bulles à l’ouverture ?
\item Les bulles prouvent-elles une réaction chimique ?
\end{questions}
\end{exercice}
\end{minipage}
\par\noindent\begin{minipage}{\linewidth}
\begin{exercice}[2]{Choisir une méthode}
\begin{questions}
\item On veut séparer l’eau et un sel dissous non volatil. Quelle méthode peut recueillir l’eau ?
\item Quelle méthode peut laisser le sel comme résidu ?
\end{questions}
\end{exercice}
\end{minipage}
\clearpage\markright{Je raisonne}
\rubrique{Je raisonne}
\par\noindent\begin{minipage}{\linewidth}
\begin{exercice}[3]{Interpréter une baisse de pH}
\begin{questions}
\item Un échantillon passe de pH 8,2 à pH 7,8. Est-il devenu acide au sens du classement scolaire ?
\item Peut-on parler d’acidification ?
\end{questions}
\end{exercice}
\end{minipage}
\par\noindent\begin{minipage}{\linewidth}
\begin{exercice}[3]{Le résultat trop précis}
\begin{questions}
\item Une bandelette présente une couleur intermédiaire entre pH 5 et 6. Un élève note 5,372. Critiquer.
\item Quelle écriture convient mieux ?
\end{questions}
\end{exercice}
\end{minipage}
\par\noindent\begin{minipage}{\linewidth}
\begin{exercice}[3]{Une comparaison équitable}
\begin{questions}
\item Deux groupes comparent la solubilité de deux sels avec des températures et volumes d’eau différents. Peut-on comparer directement les masses dissoutes ?
\item Que faudrait-il harmoniser ?
\end{questions}
\end{exercice}
\end{minipage}
\par\noindent\begin{minipage}{\linewidth}
\begin{exercice}[3]{Diluer sans deviner}
\begin{questions}
\item On ajoute de l’eau à une solution acide usuelle. Le pH se rapproche généralement de quelle valeur ?
\item Peut-on affirmer qu’il double quand le volume double ?
\end{questions}
\end{exercice}
\end{minipage}
\par\noindent\begin{minipage}{\linewidth}
\begin{exercice}[3]{L’eau prétendument potable}
\begin{questions}
\item Une publicité affirme : « Filtrée, limpide, donc potable. » Quels critères manquent ?
\item Proposer une conclusion scientifiquement correcte.
\end{questions}
\end{exercice}
\end{minipage}
\pcPage{Défi du chapitre}{Solutions, solubilité et pH}
\phantomsection\label{defi-ex-4}
\pcDefiTitre{4}{Préparer sans gaspiller}
\begin{questions}
\item Une donnée indique $30\unite g$ solubles pour $100\unite{mL}$ d’eau. On utilise $80\unite{mL}$ et $28\unite g$ de solide. Calculer la masse dissoute maximale et le dépôt.
\item Quel volume minimal d’eau permettrait théoriquement de dissoudre les $28\unite g$ ?
\item Pourquoi ce minimum doit-il être discuté expérimentalement ?
\end{questions}
\pcGrille
